% BHU PYQ -- Physics: Applied Electronics and Fiber Optics (2025-26 NEP)
% Paper Code: PHYMN31 / PHYMJ31 / PHYMV31 (Sem III Re-Mid Term / Sessional Exam, Batch P1)
\documentclass[11pt,a4paper]{article}
\usepackage[a4paper,top=2.5cm,bottom=2.5cm,left=2.8cm,right=2.8cm,headheight=14pt]{geometry}
\usepackage{amsmath,amssymb}
\usepackage[T1]{fontenc}
\usepackage[utf8]{inputenc}
\usepackage{lmodern,microtype,fancyhdr,enumitem,hyperref,lastpage,parskip}

\hypersetup{colorlinks=true,urlcolor=black,linkcolor=black,pdftitle={PHYMN31: Applied Electronics -- BHU Sem III Mid-Term (Batch P1) 2025-26 (NEP)}}

\pagestyle{fancy}\fancyhf{}
\renewcommand{\headrulewidth}{0.4pt}\renewcommand{\footrulewidth}{0.4pt}
\fancyhead[L]{\small \textit{Banaras Hindu University}}
\fancyhead[R]{\small \textit{PHYMN31: Applied Electronics (Batch P1 Mid-Term 2025-26)}}
\fancyfoot[C]{\small Page \thepage\ of \pageref{LastPage}}

\newcommand{\pts}[1]{\hfill \textit{[#1]}}

\begin{document}

\begin{center}
  {\large\bfseries BANARAS HINDU UNIVERSITY}\\[2pt]
  {\normalsize B.Sc. (Hons.) Semester III Sessional / Re-Mid Term Examination 2025-26 (NEP)}\\[6pt]
  \rule{0.8\linewidth}{0.5pt}\\[6pt]
  {\large\bfseries DEPARTMENT OF PHYSICS}\\[4pt]
  {\large\bfseries Paper Code: PHYMN31 / PHYMJ31 / PHYMV31 : Applied Electronics and Fiber Optics}\\[2pt]
  {\normalsize (Re-Mid Term Examination --- Batch P1)}\\[4pt]
  {\normalsize Time: 1 Hour \hfill Maximum Marks: 20}
\end{center}
\rule{\linewidth}{0.4pt}
\smallskip

\noindent \textit{Instructions: Answer all questions as specified. Draw neat and labelled diagrams wherever necessary.}

\bigskip

\begin{enumerate}[label=\textbf{\arabic*.},leftmargin=*,itemsep=16pt]

  \item With the help of a neat diagram, explain a cathode ray oscilloscope (CRO), and discuss the construction and working of the main components of a cathode ray tube (CRT).

  \item Draw a neat circuit diagram of a half wave rectifier. Find the efficiency and the Ripple Factor of a full wave rectifier.

  \item With the help of a neat diagram, explain the formation of the depletion layer and potential barrier in a p-n junction diode.

\end{enumerate}

\bigskip
\begin{center}
  \textbf{***}
\end{center}

\end{document}
